\documentclass[11pt,a4paper]{article}
\usepackage[T1]{fontenc}
\usepackage{lmodern,amsmath,amssymb,amsthm,mathtools}
\usepackage[margin=25mm,headheight=15pt]{geometry}
\usepackage{booktabs,array,tabularx,longtable}
\usepackage[dvipsnames]{xcolor}
\usepackage{microtype,enumitem,fancyhdr}
\usepackage[colorlinks=true,linkcolor=MidnightBlue,citecolor=MidnightBlue,urlcolor=MidnightBlue]{hyperref}
\usepackage{bookmark}
\hypersetup{pdftitle={Native Intensity Accounting and Conditional Carrier Crossings in the Tri-Octagon Model - v0.1},pdfauthor={Hilmir Frímann Halldórsson},pdfsubject={A bounded assessment of central six-opening energy transport. Author-approved repository research edition v0.1; internal mathematical acceptance, Codex manuscript and layout checks; no external peer review.}}
\setlength{\parskip}{4pt}\setlength{\parindent}{0pt}\setlength{\emergencystretch}{3em}
\setlist{itemsep=2pt,topsep=4pt}
\pagestyle{fancy}\fancyhf{}
\fancyhead[L]{\small Native intensity and conditional carrier crossings}
\fancyhead[R]{\small v0.1}\fancyfoot[C]{\thepage}
\renewcommand{\headrulewidth}{0.3pt}
\numberwithin{equation}{section}
\newtheorem{proposition}{Proposition}[section]
\newtheorem{theorem}[proposition]{Theorem}
\theoremstyle{remark}\newtheorem{remark}[proposition]{Remark}
\newcommand{\R}{\mathbb R}\newcommand{\C}{\mathbb C}\newcommand{\one}{\mathbf1}
\newcommand{\eps}{\varepsilon}\newcommand{\V}{\mathcal V}\newcommand{\D}{\mathcal D}
\newcommand{\re}{\operatorname{Re}}\newcommand{\im}{\operatorname{Im}}
\newcommand{\Argz}{\operatorname{Arg}_0}\newcommand{\diag}{\operatorname{diag}}
\newcommand{\ez}{e_z}\newcommand{\dd}{\,\mathrm d}
\newcolumntype{Y}{>{\raggedright\arraybackslash}X}
\begin{document}
\hypersetup{pageanchor=false}
\begin{titlepage}
\vspace*{6mm}
{\small\color{MidnightBlue} TRI-OCTAGON MATHEMATICAL RESEARCH\par}
\vspace{8mm}
{\LARGE\bfseries\raggedright Native Intensity Accounting and\\ Conditional Carrier Crossings\\ in the Tri-Octagon Model\par}
\vspace{6mm}
{\large A bounded assessment of central six-opening energy transport\par}
\vspace{7mm}
{\large Hilmir Fr\'imann Halld\'orsson\par}
\vspace{3mm}
{Author-approved repository research edition v0.1\quad 2 November 2026\par}
\vspace{5mm}
\begin{abstract}
We consolidate an accepted analytical assessment of whether energy leaves the Tri-Octagon's central region through its three upper and three lower openings. The unchanged native triad admits exact component and total intensity balances. An antisymmetric channel-exchange representation coexists with graph mismatch and the complete finite-step square; it is neither unique from its divergence nor a spatial opening current. We retain the exact potential and synchronization accounts and the separately proved descent theorem on its specified coherent domain. A separately adopted periodic carrier, C1, has a genuine storage, signed source-work and local flux balance, together with an explicit conditional native-readout-to-two-drive map. Its registered circle crosses one pair of finite host-notch measurement triangles. Combining the existing tangents and normals gives signed cut powers and a short-arc storage balance. These cuts are not the six central broken opening boundaries. No reciprocal native debit for carrier work or six-opening transfer allocation is established. Phase reconstruction, staged observations and retained Z examples clarify what direction and placement data do and do not identify. The result is a positive conditional balance-and-crossing statement, not a demonstration of physical emission or a nonzero flowing trajectory.
\end{abstract}
\vfill
{\small Research questions and programme: Hilmir Fr\'imann Halld\'orsson (pzychozen).
GPT provided scientific direction and project-internal acceptance of the underlying assessment.
Codex prepared this manuscript and performed fidelity, build, and final page checks.
The present typeset manuscript has not separately been inspected by GPT.
No external peer review, journal acceptance, DOI, or physical experiment is claimed.}
\end{titlepage}
\setcounter{page}{2}\hypersetup{pageanchor=true}
\section{Question, scope and native quantities}
The author asks whether energy leaves the central region through its three upper and three lower openings. The accepted assessment \cite{assessment} gives two positive results: exact accounting for native channel strengths, and a conditional carrier with signed transport through two registered host-notch crossings. It also identifies two missing interfaces: a reciprocal native debit for carrier work, and a transfer registration to the six central broken boundaries. This note presents that result without changing the model.

The author's flower-blooming analogy motivates an outward-behavior question; it does not define petals or a mechanical unfolding. The earlier assistant narrative of inward arrival, redirection, and a ``next available'' opening is an unverified interpretation. No such sequence, golden-ratio controller, material attachment, or new transfer law is assumed. The words intensity, potential, signal, and flux will retain different mathematical meanings.

\subsection{The unchanged triad}
Let $\Omega\in\C^3$, with real native coefficients $\eps,g,\lambda,k_i$, and write
\begin{equation}\label{eq:quantities}
 s_i=|\Omega_i|^2,\quad I=\sum_i s_i,\quad
 S_{ij}+iA_{ij}=\overline\Omega_i\Omega_j,\quad
 D_{ij}=s_i+s_j-2S_{ij}=|\Omega_i-\Omega_j|^2.
\end{equation}
Each undirected edge is counted once in sums over $i<j$. Put $L_N=\one\one^T-3\mathrm{Id}$. The native prestage, increment and simultaneous synchronization are \cite{nrg,potential}
\begin{align}
 b_i&=\eps(k_i-s_i)\Omega_i+g\sum_{j\ne i}(\Omega_j-\Omega_i),
 &\widetilde\Omega&=\Omega+b,\label{eq:pre}\\
 \delta_i&=\lambda\sum_{j\ne i}\sin\!\left(3[\Argz\widetilde\Omega_j-\Argz\widetilde\Omega_i]\right),
 &F_i(\Omega)&=e^{i\delta_i}\widetilde\Omega_i.\label{eq:map}
\end{align}
Every kick uses the same prestage. The convention is $\Argz(0)=0$; when $\lambda=0$, the implementation returns the prestage without extracting phases. The increment $b$ is not the complete displacement $F(\Omega)-\Omega$. No physical timestep is specified by the integer native index.

\begin{table}[htbp]\centering\small
\caption{Typed quantities. A spatial assignment does not by itself make a flux.}\label{tab:types}
\begin{tabularx}{\textwidth}{@{}p{.18\textwidth}YY@{}}\toprule
Quantity & Type & Existing meaning\\\midrule
$s_i,I$ & Nonnegative native scalars & Channel strengths and their sum; not a volume density.\\
$\V$ & Quartic native scalar & Potential on six real state coordinates; no exterior recipient.\\
$D_{ij}$ & Matched edge quadratic & Squared difference after tangent-plane transport.\\
$A_{ij},C$ & Signed pair areas; channel vector & $C=(A_{23},-A_{13},A_{12})$; separate ambient attachment.\\
$J^I_{i\to j}$ & Antisymmetric edge scalar & One algebraic channel-exchange representation per native step.\\
$e,E,j$ & C1 density, storage, local flux & Defined on a separately adopted periodic circle.\\
$(U,V_{\rm ext})$ & C6 exterior record & Retarded observation; not receiver-load energy.\\\bottomrule
\end{tabularx}
\end{table}

\section{Exact finite-step intensity accounting}
\begin{proposition}[Component balance and edge-once total]\label{prop:intensity}
For the ideal map \eqref{eq:pre}--\eqref{eq:map}, at general real coefficients,
\begin{align}
 \Delta s_i&=2\eps s_i(k_i-s_i)+2g\sum_{j\ne i}(S_{ij}-s_i)+|b_i|^2,\label{eq:component}\\
 &= -\sum_{j\ne i}J^I_{i\to j}+2\eps s_i(k_i-s_i)
       -g\sum_{j\ne i}D_{ij}+|b_i|^2,\label{eq:exchange}\\
 J^I_{i\to j}&=g(s_i-s_j)=-J^I_{j\to i},\notag\\
 \Delta I&=2\eps\sum_i s_i(k_i-s_i)-2g\sum_{i<j}D_{ij}+\|b\|^2.\label{eq:total}
\end{align}
\end{proposition}
\begin{proof}
Synchronization preserves each prestage modulus. Expanding $|\Omega_i+b_i|^2$ gives $\Delta s_i=2\re(\overline\Omega_i b_i)+|b_i|^2$ and then \eqref{eq:component}. The identity
\[
2g(S_{ij}-s_i)=g(s_j-s_i)-g(s_i+s_j-2S_{ij})
\]
gives \eqref{eq:exchange}. Summation cancels the antisymmetric terms; each undirected mismatch appears at both endpoints, giving the factor two in \eqref{eq:total}. The argument includes zero components without a phase quotient.
\end{proof}

This independently derived account in the accepted assessment agrees with the earlier diagnostic and Z-source identities \cite{nrg,zmethods,assessment}. If $h_i=(k_i-s_i)\Omega_i$, its full finite square is
\begin{align}
 \|b\|^2&=\eps^2\|h\|^2+2\eps g\re(h^\dagger L_N\Omega)
                 +g^2\|L_N\Omega\|^2,\label{eq:square}\\
 \|L_N\Omega\|^2&=3\sum_{i<j}D_{ij}.\notag
\end{align}
The mixed term has no fixed sign. For $g\ge0$ the displayed mismatch contribution is nonpositive and the selected exchange points from larger to smaller intensity. Neither fixes the sign of the complete step. For $\eps>0$, the onsite sign depends on $k_i-s_i$. No sign conclusion for arbitrary negative coefficients is implied.

\begin{remark}[Divergence does not select a unique current]\label{rem:cycle}
Adding one real number $a$ to the oriented edges $1\to2$, $2\to3$, $3\to1$, and $-a$ to their reverses, changes no node divergence. Thus \eqref{eq:exchange} is a chosen algebraic representation, not a uniquely selected physical current. Signed area is different: aligned unequal positive amplitudes have $A_{ij}=0$ but can have $J^I_{i\to j}\ne0$; equal amplitudes with nonzero relative sine have $J^I_{i\to j}=0$ but $A_{ij}\ne0$. These are comparisons of formulas, not newly evolved state examples.
\end{remark}

The graph's contribution is not exhausted by exchange. The pair mismatch and whole square must remain. Nor does weighted response geometry fill the spatial gap: at a positive fixed branch $r_i$, the inherited phase-prestage matrix has $P_{ij}=gr_j/r_i$ off diagonal, $R^2P=P^TR^2$ for $R=\diag(r_i)$, and quadratic form $g\sum_{i<j}r_ir_j(f_i-f_j)^2$. Its positive conductances for $g>0$ and the ordered complete linear response $S_\phi P$ are response facts, not six-opening currents \cite[response chapter, weighted geometry]{nrg}.

\section{Potential, synchronization and the descent domain}
Use the edge-once potential and real gradient normalization
\begin{equation}\label{eq:potential}
 \V(\Omega)=\eps\sum_i\left(\frac{s_i^2}{4}-\frac{k_i s_i}{2}\right)
       +\frac g2\sum_{i<j}D_{ij},\qquad b=-\nabla_{\R^6}\V(\Omega).
\end{equation}
Taylor's exact prestage account is
\begin{equation}\label{eq:taylor}
 \V(\Omega+b)-\V(\Omega)=-\|b\|^2+
 \int_0^1(1-t)b^T\nabla^2\V(\Omega+tb)b\dd t.
\end{equation}
For a polynomial form put $q_i=2\re(\overline\Omega_i b_i)+|b_i|^2$. Direct expansion gives
\begin{align}
 \Delta\V_{\rm pre}={}&\eps\sum_i\left[\frac{(s_i-k_i)q_i}{2}+\frac{q_i^2}{4}\right]\notag\\
 &+g\sum_{i<j}\re\!\left[\overline{(\Omega_i-\Omega_j)}(b_i-b_j)\right]
       +\frac g2\sum_{i<j}|b_i-b_j|^2.\label{eq:poly}
\end{align}
For $d_{ij}=\delta_j-\delta_i$, conjugate-pair rotation gives
\begin{align}
 S_{ij}^+&=\widetilde S_{ij}\cos d_{ij}-\widetilde A_{ij}\sin d_{ij},&
 A_{ij}^+&=\widetilde A_{ij}\cos d_{ij}+\widetilde S_{ij}\sin d_{ij},\label{eq:rotation}\\
 \Delta\V_{\rm sync}&=g\sum_{i<j}\left[\widetilde S_{ij}(1-\cos d_{ij})+
                  \widetilde A_{ij}\sin d_{ij}\right].\label{eq:sync}
\end{align}
The onsite potential is unchanged by this phase stage, while its edge term changes through $-g\sum S_{ij}$. Synchronization's zero intensity increment therefore does not mean zero potential change or no influence on later evolution.

\begin{theorem}[Inherited coherent-domain descent, not re-proved here]\label{thm:descent}
At precisely
\begin{equation}\label{eq:tuple}
 \eps=\frac1{10},\quad g=\frac1{60},\quad\lambda=\frac1{300},\quad
 k=\left(\frac{999}{1000},1,\frac{1001}{1000}\right),
\end{equation}
let $\D$ have $9/10\le|\Omega_i|\le11/10$ and ordinary phase width at most $\pi/12$, with arbitrary common orientation. The accepted theorem \cite[finite-domain descent]{potential} gives
\begin{equation}\label{eq:descent}
 \V(F\Omega)-\V(\Omega)\le-\frac{16869}{20000}\|b\|^2,\qquad\Omega\in\D.
\end{equation}
The domain is forward invariant with persistent nonzero prestages. Equality in monotonicity holds exactly on the inherited aligned fixed circle.
\end{theorem}

The width is ordinary full-turn phase width, not phase modulo $2\pi/3$; the factor multiplies the squared prestage increment. No profile is solved here. The theorem does not apply automatically to the distinct historical Z coefficient tuples in Section~\ref{sec:readouts}.

In these inspected native definitions, V and I have no accompanying exterior reservoir, six-opening boundary term, or source-work debit. Potential decrease, intensity change and emitted energy are consequently different claims. Equation~\eqref{eq:descent} supplies no destination for a decrease. This is an inventory of the existing definitions, not a theorem prohibiting a future transport extension.

\section{The existing conditional carrier and exterior observation}
\subsection{What the native-to-source interface actually supplies}
The paired-driving paper \cite{pairedII} explicitly adopts C1 and C6 as comparison constructions. They are not part of native F. For one named readout family $v_{\rm in}=M_{\rm macro}$ or raw $C$, its input chain is
\begin{equation}\label{eq:input}
 \Omega,\text{ settings}\longmapsto v_{\rm in}\longmapsto
 W_{\rm reg}=A_{\rm reg}v_{\rm in}\longmapsto
 q_\pm=n_\pm^TW_{\rm reg}\longmapsto f_\pm=\kappa_{\rm in}q_\pm,
\end{equation}
where $A_{\rm reg}=\lambda_gR_z(\delta_g)$ is the fixed registration and $\kappa_{\rm in}>0$ converts data units to source strength. It is not energy efficiency. The normals are specified in Section~\ref{sec:geometry}. With $W_{\rm reg}=W_ru+W_av+W_z e_z$,
\begin{equation}\label{eq:sectors}
 f_s=\tfrac12(f_++f_-)=\frac{2\kappa_{\rm in}}{\sqrt7}W_r,
 \qquad f_a=\tfrac12(f_+-f_-)=\frac{\sqrt3\kappa_{\rm in}}{\sqrt7}W_z.
\end{equation}
This projection loses $W_a$. Equal projected histories and equal carrier preparation give equal downstream states by the source's uniqueness theorem. Thus a genuine conditional native-to-drive map exists, but does not preserve every native observable.

The same source provides a finite common hold and end gate for a specified record:
\begin{equation}\label{eq:hold}
 f_\pm(t)=\sum_{j=0}^{N_s-1}f_{\pm,j}\mathbf1_{[j\Delta t,(j+1)\Delta t)}(t),
 \qquad \Delta t>0,
\end{equation}
with zero input outside that window. The cadence is symbolic comparison data, not a derived native physical time. No forcing record or schedule is instantiated here.

\subsection{Storage, signed work and local flux}
On an oriented circle of length $L_c$, let $s_\pm$ be its two registered contacts and $0<\eps_c<\eta$ the source's profile width. The periodic normalized tents and forcing are
\begin{equation}\label{eq:forcing}
 b_{\eps_c}(s)=\eps_c^{-1}\left(1-\frac{|\operatorname{wrap}_{L_c}(s)|}{\eps_c}\right)_+,
 \quad b_\pm(s)=b_{\eps_c}(s-s_\pm),\quad g_c=b_+f_++b_-f_-.
\end{equation}
Here $\operatorname{wrap}_{L_c}$ takes values in $[-L_c/2,L_c/2)$; the profiles have unit integral and disjoint supports. The width is arclength, not tube radius, and is unrelated to native $\eps$.

C1 is the separately adopted system
\begin{equation}\label{eq:C1}
 p_t+c(r_c)_s=g_c,\qquad (r_c)_t+c p_s=0,\qquad c>0,
\end{equation}
on $\R/L_c\mathbb Z$, with independently specified real initial data. Its variables $p,r_c$ are not native channel components. Define
\begin{equation}\label{eq:storage}
 e=\frac{p^2+r_c^2}{2},\quad j=cpr_c,\quad E=\int_0^{L_c}e\dd s,
 \quad y_\pm=\int_0^{L_c}b_\pm p\dd s.
\end{equation}
The existing well-posedness and work theorem \cite[carrier section]{pairedII} gives, for periodic $H^1$ preparation and continuous real drives, a unique solution in
$C_tH^1_s\cap C_t^1L^2_s$ on finite intervals, with
\begin{equation}\label{eq:work}
 \boxed{e_t+j_s=pg_c,\qquad \frac{dE}{dt}=f_+y_++f_-y_-.}
\end{equation}
The local identity is distributional. Multiplying the two equations by $p,r_c$ and adding gives its product rule; integration removes the periodic derivative. The source also gives a continuous $L^2$ mild solution for $L^2$ preparation and time-integrable spatial-$L^2$ forcing. In that setting the global balance is integrated in time, with derivative almost everywhere. Source-off E is conserved.

In the normalized characteristics $w_\pm=(p\pm r_c)/\sqrt2$,
\begin{equation}\label{eq:characteristic}
 e=\tfrac12(w_+^2+w_-^2),\qquad
 j=\tfrac c2(w_+^2-w_-^2),\qquad |j|\le ce.
\end{equation}
The inherited finite-time bound is
\begin{equation}\label{eq:energybound}
 \sqrt{E(t)}\le\sqrt{E(0)}+\frac1{\sqrt{3\eps_c}}
                   \int_0^t\|f(u)\|_{\R^2}\dd u.
\end{equation}
These are valid model-internal energy and transport statements without an SI calibration. If $[p]=[r_c]=\mathrm U$ and $[s]=\mathrm L$, then $[E]=\mathrm U^2\mathrm L$ and $[j]=\mathrm U^2\mathrm L/\mathrm T$. The spatially integrated directional flux $\int j\dd s$ has another length factor; it is not power through a section.

Source work is signed and depends on response as well as drive. Different allowed preparations at the same f can give different initial $f_+y_++f_-y_-$. Work is not $q_+^2+q_-^2$, $\|v_{\rm in}\|^2$, or a specified native loss. Superposition also introduces quadratic storage cross terms: the storage of $\alpha(p_1,r_1)+\beta(p_2,r_2)$ includes $\alpha\beta\int(p_1p_2+r_1r_2)\dd s$.

For zero preparation, the source's fixed-clock macro comparison gives proportional real port histories and zero spatially integrated directional bias; this does not imply zero local j. Its clock-only and limited native chiral comparisons have their own scope and are not a generated six-opening trajectory. We reuse those distinctions without rerunning them. The C1 source explicitly includes no reciprocal debit in the unchanged native recurrence \cite[carrier and conclusion]{pairedII}. An equality between native loss and integrated source work would be an additional equation with matching units and time intervals, not a consequence of \eqref{eq:input}.

\subsection{C6: exterior records without loading}
For fixed $X\notin\Gamma$, known orthonormal frame B, $d_X(s)=\|X-\Gamma(s)\|$, independent speed $c_e>0$ and gain $g_o>0$, C6 adopts
\begin{align}
 U(X,t)&=\frac{g_o}{4\pi}\int_0^{L_c}\frac{p(s,t-d_X(s)/c_e)}{d_X(s)}\dd s,\notag\\
 V_{\rm ext}(X,t)&=\frac{g_o}{4\pi}\int_0^{L_c}
       \frac{r_c(s,t-d_X(s)/c_e)\hat t(s)}{d_X(s)}\dd s,
 &Y_X&=(U,B^TV_{\rm ext}).\label{eq:C6}
\end{align}
Straight-line delay, $1/d$ weighting, tangent encoding, and absence of screening or loading are assumptions. The source proves a bounded causal history map; continuous time into $H^1$ supplies pointwise records \cite[exterior section]{pairedII}. The vector $r_c\hat t$ is not carrier flux j, raw native C, or an assigned material current.

No C6 record feeds back into \eqref{eq:C1} or \eqref{eq:work}. The source separately identifies missing reciprocal core exchange, ambient source/energy accounting, and receiver loading. A formal outgoing line-source interpretation does not provide a global near-source energy or self-force construction. Neither E nor $|U|^2+\|V_{\rm ext}\|^2$ is declared native physical energy.

\section{Central boundaries and finite host-notch patches}\label{sec:geometry}
\subsection{The actual six central objects}
Use width-one geometry with spatial center $O=(0,r_g,0)$, $r_g=\sqrt3/6$. Coordinates in this subsection are centered at O. Put
\begin{equation}\label{eq:constants}
 d_g=1-\sqrt2/2,\quad a_g=d_g/2,\quad h_g=(\sqrt2-1)/2,\quad H_0=1/2,
\end{equation}
and, in the horizontal plane,
\begin{equation}\label{eq:directions}
 e_i=(\cos(-\pi/2+i\pi/3),\sin(-\pi/2+i\pi/3)),\quad
 n_k^{\rm geom}=e_{2k},\quad g_k=e_{2k+1},\quad j_k=Jg_k,
\end{equation}
where $J(X,Y)=(-Y,X)$. For the static PLUS family,
\begin{align}
 A_\tau&=\mathrm{Id}-(\tau/r_g)J
       =\sqrt{1+(\tau/r_g)^2}R_{-\delta},
 &\delta&=\arctan(\tau/r_g),\quad0\le\tau\le h_g,\notag\\
 T_i&=(a_ge_i,h_g),&B_i&=(a_gA_\tau e_i,-h_g),\notag\\
 U_k&=(r_gn_k^{\rm geom},H_0),&L_k&=(r_gA_\tau n_k^{\rm geom},-H_0).\label{eq:vertices}
\end{align}
The central definitions \cite{central,fixedgap} give six broken two-segment boundaries
\begin{equation}\label{eq:G}
 G_k^+=U_k\,--\,T_{2k+1}\,--\,U_{k+1},\qquad
 G_k^-=L_k\,--\,B_{2k+1}\,--\,L_{k+1},\qquad k=0,1,2.
\end{equation}
The notation lists joined segments, not flow directions. Indices are cyclic. Upper label azimuths are $-30^\circ,90^\circ,210^\circ$; lower ones are shifted by $-\delta$. MINUS reflects the entire configuration by $\diag(-1,1,1)$ and retains the labels. No top-to-top closing chord, filled interior, centroid, aperture normal or hinge is assigned to G.

The refined central construction is declared additional research geometry, not a uniquely forced consequence of the folded host. Its triangulated annulus has 18 vertices, 42 edges, 24 faces and two boundary cycles. In the central slab its section complement is a connected annular region, not three isolated flow channels. Geometry alone supplies no pressure, density, velocity or boundary condition \cite{central}.

The original close sphere has radius $\sqrt{13/12-\sqrt2/2}$ about O. It is a reference, not a wall; it is distinct from an extended-line sphere and the intended tight toroid. The 180-degree support lines $O+\R g_k$ stay stationary. Geometric $\tau$, native update index, native common phase, observer clock and carrier time are different variables. Two co-present arrangements are not silently replaced by one rapidly alternating arrangement.

\subsection{A precise host-patch dictionary}
The separate finite host-notch measurement triangles are
\begin{equation}\label{eq:N}
 N_k^\sigma(\ell,\zeta)=
 \left(2r_g-\frac{\sqrt3}{2}d_g\ell\right)g_k+\zeta j_k
       +\sigma(h_g+d_g\ell)e_z,
 \quad0\le\ell\le1,\quad|\zeta|\le d_g\ell/2.
\end{equation}
Here $\ell$ is a triangle coordinate. Paper I \cite{pairedI} uses world coordinates
\begin{align}
 \Delta_j^\sigma&=\left\{O+(r_0-c_gd_gt)u_j+\xi v_j
                +\sigma(h_g+d_gt)e_z:\ 0\le t\le1,\ |\xi|\le d_gt/2\right\},\notag\\
 r_0&=1/\sqrt3,\quad c_g=\sqrt3/2,\quad
 u_j=R_z(2\pi j/3)(0,1,0),\quad v_j=e_z\times u_j.\label{eq:Delta}
\end{align}
Direct substitution of $r_0=2r_g$ and $u_0=g_1,u_1=g_2,u_2=g_0$ gives the accepted dictionary
\begin{equation}\label{eq:dictionary}
 \boxed{\Delta_j^\sigma-O=N_{(j+1)\bmod3}^\sigma.}
\end{equation}
The parameters match as $t=\ell$, $\xi=\zeta$. Two triangle sides are material rim edges; the third is a measurement-only closure, not a cap. Their oriented unit normals are
\begin{equation}\label{eq:normals}
 n_j^\sigma=\frac{2u_j+\sigma\sqrt3 e_z}{\sqrt7}.
\end{equation}
Equation~\eqref{eq:dictionary} identifies host patches across source conventions, not N with G.

The retained fixed-gap obstruction makes that distinction exact \cite[boundary obstruction]{fixedgap}. Each central broken arc contains its ring elbow with squared radius at most $15/2-21\sqrt2/4$. Every host-notch plane has squared distance from O at least $19/84+\sqrt2/14$, strictly larger. A centered rotation preserves the elbow radius and cannot put the entire arc in any host patch. This upper bound is on the elbow, not every point of the arc. Source-admitted fan-tip contacts do not remove it. It does not prohibit a different future transfer registration.

\section{Two conditional cut powers and a short-arc balance}
Paper II selects Paper I's $j=0$ circle, hence the host-notch pair $N_1^+,N_1^-$ in \eqref{eq:dictionary}. Write $u=u_0$, $v=e_z\times u$, and
\begin{equation}\label{eq:circleconstants}
 K_c=\frac{\sqrt3}{12}+\frac{\sqrt6}{4},\quad
 a_c=\frac{K_c}{1+\sqrt3/2},\quad R_c=\frac{\sqrt7a_c}{2},\quad L_c=2\pi R_c.
\end{equation}
Its adopted orientation and contact data are
\begin{align}
 \Gamma(s)&=O+K_cu+R_c(\cos\chi\,u+\sin\chi\,e_z),
 &\hat t&=-\sin\chi\,u+\cos\chi\,e_z,\quad\chi=s/R_c,\label{eq:circle}\\
 P_\pm&=O+a_cu\pm a_ce_z,&n_\pm&=(2u\pm\sqrt3e_z)/\sqrt7,\notag\\
 s_+&=L_c/2-\eta,&s_-&=L_c/2+\eta,\quad
 \eta=R_c\arctan(2/\sqrt3).\label{eq:contacts}
\end{align}
The contacts are relative-interior points of the finite patches and are not antipodal. The source's finite-crossing theorem gives
\begin{equation}\label{eq:tangents}
 \hat t(s_+)=-n_+,\qquad \hat t(s_-)=+n_-.
\end{equation}
This offset circle is an additional conditional centerline, not a material edge, not a loop through O, and not the intended tight toroid. Each candidate circle meets only its own pair of finite patches, avoiding the other four. The three geometric candidates are mutually disjoint; no planar circle contains all six selected contacts. These facts do not instantiate three carriers or their coupling \cite[normal-crossing circle and clearance]{pairedI}.

\begin{proposition}[Conditional signed cut powers and arc balance]\label{prop:crossing}
For the already adopted C1 state at spatial $H^1$ regularity, the powers through the two registered cuts, with the outward host-plane signs, are
\begin{equation}\label{eq:cuts}
 \boxed{P_{\rm out,+}(t)=-j(s_+,t),\qquad P_{\rm out,-}(t)=+j(s_-,t).}
\end{equation}
On the short arc, the integrated-time identity is
\begin{align}
 E_{\rm arc}(t_1)-E_{\rm arc}(t_0)
  &=\int_{t_0}^{t_1}\left[\int_{s_+}^{s_-}pg_c\dd s
            -P_{\rm out,+}-P_{\rm out,-}\right]\dd t,\label{eq:arcint}\\
 E_{\rm arc}(t)&=\int_{s_+}^{s_-}e(s,t)\dd s.\notag
\end{align}
Its derivative form holds almost everywhere:
\begin{equation}\label{eq:arc}
 \frac{dE_{\rm arc}}{dt}=\int_{s_+}^{s_-}pg_c\dd s+j(s_+,t)-j(s_-,t).
\end{equation}
\end{proposition}
\begin{proof}
The source normals point away from the O-side half-spaces because
$n_\pm\cdot(P_\pm-O)=a_c(2+\sqrt3)/\sqrt7>0$.
The one-dimensional cut power is $j\hat t\cdot n_\pm$, so \eqref{eq:tangents} gives \eqref{eq:cuts}. In one spatial dimension $H^1$ supplies continuous point traces and the weak product rule used in \eqref{eq:work}. Integrating that local balance on $[s_+,s_-]$ and then in time yields \eqref{eq:arcint} and its almost-everywhere derivative. This is the accepted composition of the existing definitions \cite[Section B.4]{assessment}, not a new wave solution.
\end{proof}

The source tents straddle their contact cuts. Therefore the work restricted to the arc is generally not the entire $f_+y_++f_-y_-$. The powers in \eqref{eq:cuts} are one-dimensional cut powers, not surface-flux integrals through filled central apertures. Point-flux statements here are not extended to arbitrary spatial $L^2$ preparation; the weaker source theorem supplies the global integrated balance.

Positive local j is incoming at the upper contact and outgoing at the lower; negative j reverses those signs locally. The two fluxes need not have equal value or sign at a fixed time, and no permanent positive sign is asserted. The carrier is periodic: characteristics continue around it, rather than escaping or reflecting at the source ports. ``Inner'' and ``outer'' midpoint labels describe locations, not an inward-then-outward routing law from O.

Even contact is weaker than crossing. Paper I's frozen unit macro reference
$\mathcal M_j(\nu)=O+a_c\cos3\nu(\cos\nu\,u_j+\sin\nu\,v_j+e_z)$ meets $P_j^+$ at zero, but has tangent $a_cv_j$, tangent to the patch plane. The selected circle instead crosses normally. That reference curve is not a generated native history. No shared path or ``next available'' central opening follows \cite[shared-point distinction]{pairedI}.

\section{Readouts, phase information and retained observations}\label{sec:readouts}
\subsection{Center, based vectors and matched transport}
The complex mean $m=\sum_i\Omega_i/3$ is not the spatial center O. Although $(L_N\Omega)_i=3(m-\Omega_i)$, that is a channel-plane correction. The exact mean-step identity is \cite[Section 2]{relaxation}
\begin{equation}\label{eq:mean}
 \Delta m=\frac\eps3\sum_i\Omega_i(k_i-s_i)
        +\frac13\sum_i\widetilde\Omega_i(e^{i\delta_i}-1).
\end{equation}
Only the direct graph sum cancels. With synchronization off, graph-induced state changes can affect later nonlinear mean steps. With synchronization on, the second term can depend on g through the prestage already in the first full step. There is no implied inward spatial path toward O.

The adopted decoder \cite{nrg,phase} gives a based face vector and endpoint
\begin{equation}\label{eq:decoder}
 v_i=D_i\Omega_i=\re(\Omega_i)t_i+\im(\Omega_i)e_z,
 \qquad(c_i,v_i),\qquad c_i+v_i,
\end{equation}
where $t_i,e_z$ are an orthonormal frame in the face tangent plane and $c_i$ is its fixed center. Its norm is $\sqrt{s_i}$, not a travel distance; the endpoint need not lie on the finite face and is not a particle-position law. Matched transport is
\begin{equation}\label{eq:matched}
 T_{ij}=t_it_j^T+e_ze_z^T,\qquad D_{ij}=\|v_i-T_{ij}v_j\|^2.
\end{equation}
It is an isometry between the tangent planes, rank two as an ambient map, not a flow along material seams.

For the accepted exact frames, $n_i=-g_{i-1}$ and $c_i-O=r_gn_i$. Thus the positioned unoriented lines agree:
\begin{equation}\label{eq:axes}
 c_i+\R n_i=O+\R g_{i-1},\qquad
 D_i(e^{i\chi}\omega)=R_{n_i}(\chi)D_i\omega=R_{g_{i-1}}(-\chi)D_i\omega.
\end{equation}
Pure phase rotates each decoded vector about its stationary support line with the directed sign shown. The full F also changes amplitudes and relative phases. The published obstruction to one fixed-label proper rotation working for every admitted triple concerns a different quantified question; it does not prohibit outward transport or prove emission \cite{phase}.

Raw chirality is $C=\re\Omega\times\im\Omega=(A_{23},-A_{13},A_{12})$ in channel coordinates. Its ambient horizontal axial attachment is $W=T_fC$, where $T_f=[t_1\ t_2\ t_3]$; the scalar $\one^TC$ supplies a distinct vertical polar component. Raw C does not obey the same ambient vector transformation law as W \cite[state--shell attachment and observables]{nrg}.

\subsection{Phase reconstruction versus staged Z}
On the exact domain of Theorem~\ref{thm:descent}, the five retained coordinates
$y=(s_1,s_2,s_3,A_{12},A_{13})$ determine full coherence. To display the positive-root convention, put
\begin{align}
 p_{12}&=\sqrt{s_1s_2-A_{12}^2},&p_{13}&=\sqrt{s_1s_3-A_{13}^2},\notag\\
 z(y)&=\left(\sqrt{s_1},\frac{p_{12}+iA_{12}}{\sqrt{s_1}},
                         \frac{p_{13}+iA_{13}}{\sqrt{s_1}}\right)^T.\label{eq:representative}
\end{align}
The narrow sector fixes the positive roots. The accepted reconstruction \cite{potential,phase} supplies a closed map G and unit multiplier $\kappa(y)=F_1(z(y))/|F_1(z(y))|$ such that
\begin{equation}\label{eq:reconstruction}
 \Omega_n=u_0\left(\prod_{j=0}^{n-1}\kappa(G^jy_0)\right)z(G^ny_0),
 \qquad u_0=\Omega_{1,0}/|\Omega_{1,0}|.
\end{equation}
Persistent nonzero prestages justify this formula. No iteration is computed here. The original common phase $u_0$ is necessary when only y is retained.

The separately defined staged readout and placement are \cite{zmethods}
\begin{align}
 K&=\lambda_{vp}\frac{\sqrt I}{1+\sqrt I},&E_o&=e^{-\gamma t_o},\notag\\
 H_\vartheta&=\cos[3(\vartheta-\ell_o)](\cos\vartheta,\sin\vartheta,1),
 &\vartheta&=2\pi(q\bmod N)/N,\notag\\
 M_{\rm macro}&=KE_oH_\vartheta,&Z&=\alpha M_{\rm macro}+\beta C,
 \qquad X=O+\mu R_z(\psi)Z.\label{eq:Z}
\end{align}
The scale, yaw, clock and configuration are external inputs; raw C here is not replaced by W. At fixed settings Z can change with native I and C, but discards common phase. Two states with the same retained y and different $u_0$ have identical retained and fixed-setting staged histories, while their nonzero decoded vectors rotate in their own planes. Consequently a genuinely common-phase-dependent outlet quantity cannot be a function of that staged history alone. This conditional identifiability statement neither requires all outlet rules to depend on phase nor denies native evolution. Adding $u_0$ recovers decoded vectors, not an unprovided transfer law.

EMA has a separate information path. Its innovation uses
$J_{\rm cubic}=\im(\Omega_1\overline\Omega_2\Omega_3)
=\im[u z_1\overline z_2z_3]$, normalized by $1+|J_{\rm cubic}|$, together with prior memory and the implementation's literal $0.01$ innovation coefficient. An observation from current memory instead uses current I, C, that memory and the external settings; it neither evaluates the cubic nor advances memory. The prescribed native-update/memory-update order is retained. Neither observer nor memory is advanced in this build \cite{phase,zmethods}.

\subsection{What the saved Z examples show}
Z01 uses two saved binary64 states at $(\eps,g,\lambda)=(0.05,0.2,0.001)$, $k=(1,1,1)$. Z02 and Z03 share five saved states from the earlier 120-digit producer at $(0.1,0.2,0.17)$, $k=(2,3,4)$. They are distinct constructed histories, not representative sampling; neither is covered automatically by Theorem~\ref{thm:descent}. The following are retained observations, not new evaluations \cite{zmethods,assessment}.

Z01's frozen nearest label is $G_0^+$ at both states, while the second scheduled observation has label $G_1^+$. In Z02 the five frozen labels are all $G_0^+$; the scheduled labels are $G_0^+,G_0^+,G_2^-,G_0^-,G_1^+$. In its recorded replacement order, all three switches occur at clock replacement. This finite-window result is not a general theorem that label changes are observer-made.

All primary Z01/Z02 placed points are inside the original sphere and off both G and N. The source-defined forward ray $O+t(X-O)$, $t\ge0$, has no finite-notch hit in Z01; Z02 reports one for frozen F1 at $N_0^+$, though F1 itself is not on the patch. A nearest label is neither incidence nor passage, and ray incidence is not flux. No face-based vector has its origin moved to O to make this test.

Z03's native-source budgets account for all four frozen displacements with its retained numerical residual qualifications. On $2\to3$, $\Delta I\simeq+0.00754918$ accompanies $\Delta\text{radius}\simeq-0.000590475$; the full finite square is essential to the positive intensity sign. On $3\to4$, radius increases by approximately $0.00780087$ while height decreases by approximately $0.00120391$. This is a genuine outward radial endpoint change of the defined readout, not an observation of energy leaving a region. Signed contributions can cancel. Comparison replacements in the native/envelope/clock telescope are not intermediate native times, continuous paths, passage events or speeds.

\section{What is established and what remains unassigned}
The existing model has exact native accounting and separately valid conditional transport mathematics. Equations~\eqref{eq:component}--\eqref{eq:total} account for channel changes without assigning their destination. Equations~\eqref{eq:work} and \eqref{eq:cuts}--\eqref{eq:arc} provide signed carrier work and two registered host-notch cut powers. The latter powers are conditional expressions whose signs depend on the carrier state, not evidence that a nonzero flowing trajectory or physical emission has been observed.

\begin{table}[htbp]\centering\small
\caption{Dependencies for the central six-opening question. ``Absent'' means absent from the inspected definitions, not universally impossible.}\label{tab:interfaces}
\begin{tabularx}{\textwidth}{@{}>{\raggedright\arraybackslash}p{.32\textwidth}>{\raggedright\arraybackslash}p{.18\textwidth}Y@{}}\toprule
Link & Status & Content or missing relationship\\\midrule
Native state to intensity changes & Exact algebra & Full exchange, mismatch, onsite and finite-square account.\\
Native state to potential loss & Inherited theorem & Exact tuple/domain; no exterior recipient.\\
Retained history to face vectors & With original phase & Reconstruction needs $u_0$; vectors are not flux.\\
Native state/settings to staged Z & Defined readout & Placement, nearest label and ray predicates are not passage.\\
Staged history to a phase-sensitive outlet & Underdetermined & Common phase is discarded; restoring it still does not define transfer.\\
Named readout to two C1 drives & Conditional map & Registration/gain and comparison-time bridge exist.\\
C1 state to signed host-cut power & Conditional exact result & Selected $N_1^\pm$ pair, prescribed normals, adequate traces.\\
Native loss to C1 work & Absent & Reciprocal native amount/work debit is required.\\
Host cuts to six central G boundaries & Absent & A declared transfer registration or direct six-transfer observable is required.\\
C1 history to C6 records & Conditional map & Ambient energy account and receiving load remain separate.\\\bottomrule
\end{tabularx}
\end{table}

Even hypothetically equating a loss with six nonnegative amounts would only impose a total. For $Q>0$, both $(Q,0,0,0,0,0)$ and $(0,Q,0,0,0,0)$ sum to Q; an interior allocation has five independent scalar freedoms. For $Q=0$, nonnegativity forces all amounts to vanish. No equation identifying Q with native loss is adopted here.

To state which central opening receives an amount leaving the core requires a named balance quantity and an explicit transfer observable attached to the six G objects. A surface-flux interpretation additionally needs declared oriented apertures/control surfaces, which the broken arcs do not supply. Using the existing C1 route instead requires a reciprocal native debit and a registration from its host-notch cuts to G. These are bounded open modeling decisions, not a proposal to implement them in this edition. Six-central-opening emission remains unestablished.

\section{Reproducibility, evidence and internal review}
This is an editorial consolidation of the accepted assessment \cite{assessment}, not a new scientific investigation. The prior standalone symbolic record reports 29 exact residuals in nine groups; it is retained as prior evidence and was not rerun. No native trajectory, state sample through F, profile solve, carrier wave solution, new geometry, observer/EMA advance or historical checker was executed. Existing numerical evidence remains qualified by its original producer and residual records, not authenticated anew by typesetting.

The accompanying claim/source ledger maps propositions, conditional statements and quoted numbers to the accepted proof's sections and equation labels and to exact earlier source files. The provenance inventory records full-file hashes and immutable reading references. The accepted proof is retained byte-identically as historical evidence; any local paths in it are historical locators, not build dependencies. Full earlier publications remain the primary references; no excerpt is presented as a complete publication. The manuscript requires only its local TeX sources and the existing cached TeX environment.

The scientific acceptance is project-internal and is recorded in the author's publication order. The separate GPT disposition file was not available in the bounded local lookup; its text is not reconstructed or claimed to have been read. The order expressly permits proceeding from the original proof and ledger and records the operative acceptance and restrictions. Codex performed manuscript fidelity and final document checks; the typeset manuscript has not separately been reviewed by GPT. No external peer review or blanket manuscript license is implied by public repository visibility.

\appendix
\section{Notation and source crosswalk}
\begin{table}[htbp]\centering\small
\caption{Symbols whose distinctions matter for the result.}
\begin{tabularx}{\textwidth}{@{}p{.27\textwidth}Y@{}}\toprule
Notation & Meaning and distinction\\\midrule
$b,\widetilde\Omega,F\Omega$ & Prestage increment, prestage, complete synchronized state.\\
$J^I_{i\to j},A_{ij},j$ & Channel exchange, signed area, C1 local flux: different types.\\
$L_N,L_c$ & Native graph matrix versus carrier circumference.\\
$g,g_k,g_c$ & Native coefficient, fixed support direction, carrier forcing.\\
$\eps,\eps_c$ & Native onsite coefficient versus carrier source width.\\
$r_g,r_c,R_c$ & Geometric inradius, carrier variable, selected circle radius.\\
$n_k^{\rm geom},n_i,n_j^\sigma$ & Central hexagon direction, face normal, host-notch normal.\\
$C,W,W_{\rm reg}$ & Raw chirality, ambient axial attachment, registered input vector.\\
$G_k^\sigma,N_k^\sigma,\Delta_j^\sigma$ & Central broken boundary, centered host patch, world host patch; $\Delta_j^\sigma-O=N_{j+1}^\sigma$ modulo three.\\
$\tau,n,u_0,(q,t_o),t$ & Geometry parameter, native index, initial common phase, observer clock, carrier time.\\\bottomrule
\end{tabularx}
\end{table}
The paper's intensity and potential derivations correspond to assessment A.2--A.3; C1/C6 to A.4--A.5; central geometry and readouts to B.1--B.3; and the host dictionary and cut/arc balance to B.4. Table~\ref{tab:interfaces} consolidates B.5. The expanded ledger retains original line and equation locators rather than treating this editorial section numbering as new scientific authority.

\begin{thebibliography}{99}\small\raggedright
\bibitem{assessment} Tri-Octagon research programme.
\emph{Native central energy transport: proof review}. Task NATIVE\_CENTRAL\_ENERGY\_TRANSPORT\_01, 2 November 2026.
\href{evidence/accepted_proof.md}{Byte-identical accepted proof}; prior symbolic record and original claim ledger accompany this edition. Project-internal acceptance is recorded in the publication order; see provenance for identities and review availability.

\bibitem{potential} H. F. Halld\'orsson.
\emph{Native Potential Descent and Closed Coherence Coordinates in the Tri-Octagon Map}. Research edition v0.1, 2 November 2026.
\href{https://github.com/pzychozen/trioctagon-physics/tree/634ff13b666604d5caed11cc3731dcb5a863ff6c/papers/NATIVE_POTENTIAL_AND_CLOSED_COHERENCE/publication/v0.1}{Immutable publication}. Setting, finite-domain descent, coordinates and retention sections.

\bibitem{phase} H. F. Halld\'orsson.
\emph{Phase Reconstruction and Stationary Support Axes in the Tri-Octagon Map}. Research edition v0.1, 2 November 2026.
\href{https://github.com/pzychozen/trioctagon-physics/tree/c881b682ef6a8128dd1dc316d1850b9aeb1e1c8b/papers/PHASE_RECONSTRUCTION_AND_STATIONARY_SUPPORT_AXES/publication/v0.1}{Immutable publication}. Reconstruction, spatial, rotation and observations sections.

\bibitem{nrg} Tri-Octagon research programme.
\emph{Native Response Geometry and Spectral Circulation in the Tri-Octagon Map}. Research publication v0.1, 6 October 2026.
\href{https://github.com/pzychozen/trioctagon-physics/tree/c881b682ef6a8128dd1dc316d1850b9aeb1e1c8b/papers/NATIVE_RESPONSE_GEOMETRY/manuscript}{Immutable reading sources}. Response chapter and observables/closure chapter.

\bibitem{relaxation} H. F. Halld\'orsson.
\emph{Transient Collective Turning and Coherent Relaxation in the Native Triad}. Research edition v0.1, 28 October 2026.
\href{https://github.com/pzychozen/trioctagon-physics/blob/c881b682ef6a8128dd1dc316d1850b9aeb1e1c8b/papers/COLLECTIVE_TURNING_AND_COHERENT_RELAXATION/v0.1/main.tex}{Immutable reading source}. Exact mean-step accounting.

\bibitem{central} H. F. Halld\'orsson.
\emph{Central Geometry, Z Readouts, and Chiral Response}. Author-authorized repository edition v0.3.
\href{https://github.com/pzychozen/trioctagon-physics/tree/c881b682ef6a8128dd1dc316d1850b9aeb1e1c8b/papers/CENTRAL_GEOMETRY_Z_CHIRAL_RESPONSE/publication/v0.3}{Immutable publication}. Central family, openings, section complement and placement conventions.

\bibitem{fixedgap} Tri-Octagon research programme.
\emph{NCM-FIXED-GAP-01 v0.1}. Accepted project geometry report, 31 October 2026.
\href{https://github.com/pzychozen/trioctagon-physics/blob/c881b682ef6a8128dd1dc316d1850b9aeb1e1c8b/papers/PHASE_RECONSTRUCTION_AND_STATIONARY_SUPPORT_AXES/publication/v0.1/provenance/primary/fixed_gap_report.md}{Retained report}, with its accepted review in the same primary-source directory. Opening definitions and ring-elbow obstruction.

\bibitem{zmethods} H. F. Halld\'orsson.
\emph{Understanding Z Motion in the Tri-Octagon Model: Native changes, observer effects, and fixed-geometry source accounting}. Methods note v0.1, 1 November 2026.
\href{https://github.com/pzychozen/trioctagon-physics/tree/c881b682ef6a8128dd1dc316d1850b9aeb1e1c8b/docs/methods/Z_MOTION/v0.1}{Immutable methods edition}, including retained Z01--Z03 reports. No history rerun here.

\bibitem{pairedI} H. F. Halld\'orsson.
\emph{Paired-Gap Geometry and Toroidal Observation Maps for the Folded Tri-Octagon}. Paper I, research edition v0.1.1, manuscript dated 18 October 2026.
\href{https://github.com/pzychozen/trioctagon-physics/tree/c881b682ef6a8128dd1dc316d1850b9aeb1e1c8b/SRG_REDISCOVERY/PAPERS/PAIRED_GAP_GEOMETRY_OBSERVATION/publication/v0.1.1}{Immutable publication}. Finite triangles, normal-crossing circle, clearance and shared-point distinction.

\bibitem{pairedII} H. F. Halld\'orsson.
\emph{Paired Driving and Exterior Identifiability on a Conditional Periodic Carrier}. Paper II, research edition v0.1, manuscript dated 19 October 2026.
\href{https://github.com/pzychozen/trioctagon-physics/tree/c881b682ef6a8128dd1dc316d1850b9aeb1e1c8b/SRG_REDISCOVERY/PAPERS/PAIRED_DRIVING_EXTERIOR_IDENTIFIABILITY/publication/v0.1}{Immutable publication}. Foundations, C1 carrier, native time bridge, C6 exterior and conclusion sections.
\end{thebibliography}

\end{document}
